\subsubsection{Trajectory Contribution Analysis}
\begin{enumerate}
    \item \textbf{Fortran Core Contribution Analysis Functions}
          
    \item \textbf{Trajectory Contribution Calculation}
          
          Computes integrated trajectory-level contribution between two vectors using
          cosine similarity.
          
          \textbf{Arguments:}
          \begin{itemize}
            \item \texttt{factor}: 1D numpy array of shape \texttt{(n\_timepoints,)}
                  — independent variable
                  
            \item \texttt{dependent}: 1D numpy array of shape \texttt{(n\_timepoints,)}
                  — dependent variable
                  
            \item \texttt{mode}: Processing mode (1 = cosine similarity, 2 = angle/acos)
          \end{itemize}
          \textbf{Returns:}
          \begin{itemize}
            \item \texttt{contribution}: Scalar contribution value (float)
          \end{itemize}
          \begin{lstlisting}[language=Python, caption={Trajectory-Level Contribution Calculation}]
contribution = tox_trajectory_contribution(factor, dependent, mode)
          \end{lstlisting}
          
    \item \textbf{Spike Contribution Calculation}
          
          Computes timepoint-wise spike contributions between two vectors.
          
          \textbf{Arguments:}
          \begin{itemize}
            \item \texttt{factor}: 1D numpy array of shape \texttt{(n\_timepoints,)}
                  — independent variable
                  
            \item \texttt{dependent}: 1D numpy array of shape \texttt{(n\_timepoints,)}
                  — dependent variable
                  
            \item \texttt{mode}: Processing mode (1 = Normal, 2 = RAP)
          \end{itemize}
          \textbf{Returns:}
          \begin{itemize}
            \item \texttt{contribution}: 1D numpy array of shape \texttt{(n\_timepoints,)}
                  with contribution values per timepoint (read-only)
          \end{itemize}
          \begin{lstlisting}[language=Python, caption={Spike-Wise Contribution Calculation}]
spike_contributions = tox_spike_contribution(factor, dependent, mode)
          \end{lstlisting}
          
    \item \textbf{Calculate Contributions (Allocation Version)}
          
          Calculates both spike and integrated contributions for a specific factor with
          internal memory allocation.
          
          \textbf{Arguments:}
          \begin{itemize}
            \item \texttt{trajectories}: 3D numpy array of shape \texttt{(n\_factors,
                  n\_samples, n\_timepoints)} in Fortran order
            
            \item \texttt{i\_factor}: 0-based index of the independent variable
                  
            \item \texttt{dependent\_idx}: 0-based index of the dependent variable
                  
            \item \texttt{mode}: Processing mode (1 = Normal, 2 = RAP)
          \end{itemize}
          \textbf{Returns:}
          \begin{itemize}
            \item Dictionary containing:
                  \begin{itemize}
                    \item \texttt{spikes}: Spike contributions \texttt{[n\_timepoints, n\_samples]}
                          (read-only)
                          
                    \item \texttt{trajectory}: Integrated contributions \texttt{[n\_samples]}
                          (read-only)
                  \end{itemize}
          \end{itemize}
          \begin{lstlisting}[language=Python, caption={Contributions Calculation with Internal Allocation}]
result = tox_calc_contributions(trajectories, i_factor, 
                                dependent_idx, mode)

# Access results
spike_contribs = result['spikes']        # [n_timepoints, n_samples]
traj_contribs = result['trajectory']     # [n_samples]
          \end{lstlisting}
          
    \item \textbf{Calculate Contributions (Expert Version)}
          
          Calculates both spike and integrated contributions using caller-provided work
          arrays for optimal performance.
          
          \textbf{Arguments:}
          \begin{itemize}
            \item \texttt{trajectories}: 3D numpy array of shape \texttt{(n\_factors,
                  n\_samples, n\_timepoints)} in Fortran order
            
            \item \texttt{i\_factor}: 0-based index of the independent variable
                  
            \item \texttt{dependent\_idx}: 0-based index of the dependent variable
                  
            \item \texttt{mode}: Processing mode (1 = Normal, 2 = RAP)
                  
            \item \texttt{temp\_factor\_vector}: 1D work array of shape \texttt{(n\_timepoints,)}
                  
            \item \texttt{temp\_dependent\_vector}: 1D work array of shape \texttt{(n\_timepoints,)}
          \end{itemize}
          \textbf{Returns:}
          \begin{itemize}
            \item Dictionary containing:
                  \begin{itemize}
                    \item \texttt{spikes}: Spike contributions \texttt{[n\_timepoints, n\_samples]}
                          (read-only)
                          
                    \item \texttt{trajectory}: Integrated contributions \texttt{[n\_samples]}
                          (read-only)
                  \end{itemize}
          \end{itemize}
          \begin{lstlisting}[language=Python, caption={Expert Contributions Calculation with Work Arrays}]
result = tox_calc_contributions_expert(
    trajectories, i_factor, dependent_idx, mode, 
    temp_factor_vector, temp_dependent_vector
)

# Access results (same structure as tox_calc_contributions)
spike_contribs = result['spikes']
traj_contribs = result['trajectory']
          \end{lstlisting}
\end{enumerate}

\textbf{Processing Modes:}
\begin{itemize}
    \item \texttt{mode = 1} (Normal): Returns cosine similarity for normal vectors
          
    \item \texttt{mode = 2} (RAP): Returns arccos of cosine similarity for RAP projected
          vectors
\end{itemize}

\textbf{Key Features of Python Interface:}
\begin{itemize}
    \item \textbf{Automatic Memory Management}: Handles Fortran memory layout and
          data type conversion automatically
          
    \item \textbf{Comprehensive Error Handling}: Converts Fortran error codes to
          Python exceptions
          
    \item \textbf{Memory Efficient}: Output arrays are marked as read-only to prevent
          accidental modification
          
    \item \textbf{Flexible Input}: Accepts both lists and numpy arrays with automatic
          conversion
          
    \item \textbf{Performance Optimized}: Expert versions allow pre-allocated work
          arrays for critical code paths
\end{itemize}

\textbf{Usage Example:}
\begin{lstlisting}[language=Python, caption={Complete Contribution Analysis Workflow}]
import numpy as np

# Example data
n_factors, n_samples, n_timepoints = 5, 100, 50
trajectories = np.random.rand(n_factors, n_samples, n_timepoints).astype(np.float64)

# Single vector analysis
factor_vec = np.random.rand(n_timepoints)
dependent_vec = np.random.rand(n_timepoints)

# Compute trajectory contribution
traj_contrib = tox_trajectory_contribution(factor_vec, dependent_vec, 
    mode=1)

# Compute spike contributions  
spike_contribs = tox_spike_contribution(factor_vec, dependent_vec, mode=1)

# Full trajectory analysis
result = tox_calc_contributions(trajectories, i_factor=0, 
    dependent_idx=1, mode=1)
spikes = result['spikes']        # [50, 100]
trajectory = result['trajectory'] # [100]

# Expert version with pre-allocated work arrays
temp_factor = np.empty(n_timepoints, dtype=np.float64)
temp_dependent = np.empty(n_timepoints, dtype=np.float64)
result_expert = tox_calc_contributions_expert(
    trajectories, i_factor=0, dependent_idx=1, mode=1,
    temp_factor, temp_dependent
)
\end{lstlisting}
\begin{enumerate}
    \item \textbf{Python Interface for Outlier Detection}
          
          The Python interface provides wrappers for the Fortran outlier detection routines,
          handling data type conversion, memory layout, and error checking
          automatically.
          
          \begin{enumerate}
            \item \textbf{Calculate Integrated Threshold (Expert Version)}
                  
                  Calculate empirical threshold for integrated contributions using pre-computed
                  permutations for optimal performance.
                  
                  \textbf{Arguments:}
                  \begin{itemize}
                    \item \texttt{contributions}: 1D numpy array of integrated contributions
                          [n\_samples] in Fortran order
                          
                    \item \texttt{percentile\_val}: Percentile value for threshold (0.0--100.0)
                          
                    \item \texttt{permutation}: Pre-computed permutation indices for sorted
                          contributions [n\_samples] in Fortran order
                  \end{itemize}
                  \textbf{Returns:}
                  \begin{itemize}
                    \item \texttt{threshold}: Scalar threshold value
                  \end{itemize}
                  \begin{lstlisting}[language=Python, caption={Integrated Threshold Calculation (Expert)}]
threshold = calc_integrated_threshold_expert(
    contributions, percentile_val, permutation
)
                  \end{lstlisting}
                  
            \item \textbf{Calculate Integrated Threshold (Convenience Version)}
                  
                  Calculate empirical threshold for integrated contributions with internal
                  allocation and sorting.
                  
                  \textbf{Arguments:}
                  \begin{itemize}
                    \item \texttt{contributions}: 1D numpy array of integrated contributions
                          [n\_samples] in Fortran order
                          
                    \item \texttt{percentile\_val}: Percentile value for threshold (0.0--100.0)
                  \end{itemize}
                  \textbf{Returns:}
                  \begin{itemize}
                    \item \texttt{threshold}: Scalar threshold value
                  \end{itemize}
                  \begin{lstlisting}[language=Python, caption={Integrated Threshold Calculation}]
threshold = calc_integrated_threshold(contributions, percentile_val)
                  \end{lstlisting}
                  
            \item \textbf{Detect Integrated Outliers}
                  
                  Identify outlier samples where integrated contributions exceed the empirical
                  threshold.
                  
                  \textbf{Arguments:}
                  \begin{itemize}
                    \item \texttt{contributions}: 1D numpy array of integrated contributions
                          [n\_samples] in Fortran order
                          
                    \item \texttt{threshold}: Scalar threshold value for outlier detection
                  \end{itemize}
                  \textbf{Returns:}
                  \begin{itemize}
                    \item \texttt{outlier\_mask}: 1D boolean array indicating outlier samples
                          [n\_samples] (read-only)
                  \end{itemize}
                  \begin{lstlisting}[language=Python, caption={Integrated Outlier Detection}]
outlier_mask = detect_outliers_integrated(contributions, threshold)
                  \end{lstlisting}
                  
            \item \textbf{Calculate Spike Thresholds (Expert Version)}
                  
                  Calculate empirical thresholds for spike contributions using pre-sorted
                  permutations for optimal performance.
                  
                  \textbf{Arguments:}
                  \begin{itemize}
                    \item \texttt{spike\_contribs}: 2D numpy array of spike contributions
                          [n\_timepoints, n\_samples] in Fortran order (rows=timepoints,
                          columns=samples)
                          
                    \item \texttt{percentile\_val}: Percentile value for threshold (0.0--100.0)
                          
                    \item \texttt{permutation}: Pre-computed permutation indices [n\_samples,
                    n\_timepoints] in Fortran order (each column contains sorted indices
                    for a timepoint)
                  \end{itemize}
                  \textbf{Returns:}
                  \begin{itemize}
                    \item \texttt{thresholds}: 1D numpy array of thresholds for each timepoint
                          [n\_timepoints] (read-only)
                  \end{itemize}
                  \begin{lstlisting}[language=Python, caption={Spike Thresholds Calculation (Expert)}]
thresholds = calc_spike_thresholds_expert(
    spike_contribs, percentile_val, permutation
)
                  \end{lstlisting}
                  
            \item \textbf{Calculate Spike Thresholds (Convenience Version)}
                  
                  Calculate empirical thresholds for spike contributions with internal allocations
                  and sorting.
                  
                  \textbf{Arguments:}
                  \begin{itemize}
                    \item \texttt{spike\_contribs}: 2D numpy array of spike contributions
                          [n\_timepoints, n\_samples] in Fortran order (rows=timepoints,
                          columns=samples)
                          
                    \item \texttt{percentile\_val}: Percentile value for threshold (0.0--100.0)
                  \end{itemize}
                  \textbf{Returns:}
                  \begin{itemize}
                    \item \texttt{thresholds}: 1D numpy array of thresholds for each timepoint
                          [n\_timepoints] (read-only)
                  \end{itemize}
                  \begin{lstlisting}[language=Python, caption={Spike Thresholds Calculation}]
thresholds = calc_spike_thresholds(spike_contribs, percentile_val)
                  \end{lstlisting}
                  
            \item \textbf{Detect Spike Outliers}
                  
                  Identify outliers in spike contributions by comparing against timepoint-specific
                  thresholds.
                  
                  \textbf{Arguments:}
                  \begin{itemize}
                    \item \texttt{spike\_contribs}: 2D numpy array of spike contributions
                          [n\_timepoints, n\_samples] in Fortran order (rows=timepoints,
                          columns=samples)
                          
                    \item \texttt{thresholds}: 1D numpy array of thresholds for each timepoint
                          [n\_timepoints] in Fortran order
                  \end{itemize}
                  \textbf{Returns:}
                  \begin{itemize}
                    \item \texttt{outlier\_mask}: 2D boolean array indicating outliers [n\_timepoints,
                    n\_samples] (read-only)
                  \end{itemize}
                  \begin{lstlisting}[language=Python, caption={Spike Outlier Detection}]
outlier_mask = detect_outliers_spike(spike_contribs, thresholds)
                  \end{lstlisting}
          \end{enumerate}
          
          \textbf{Key Features of the Python Interface:}
          \begin{itemize}
            \item \textbf{Automatic Type Conversion}: Input arrays automatically converted
                  to \texttt{np.float64} and \texttt{np.int32} with Fortran memory layout
                  
            \item \textbf{Dimension Validation}: Input dimensions are validated before
                  calling Fortran routines
                  
            \item \textbf{Error Handling}: Fortran error codes are converted to Python
                  exceptions via \texttt{check\_err\_code()}
                  
            \item \textbf{Memory Efficiency}: Output arrays are marked as read-only to
                  prevent accidental modification
                  
            \item \textbf{Boolean Conversion}: C integer arrays (0/1) are automatically
                  converted to Python boolean arrays
          \end{itemize}
          
          \textbf{Usage Notes:}
          \begin{itemize}
            \item Use \textbf{expert} versions when you have pre-computed
                  permutations for better performance
                  
            \item Use \textbf{convenience} versions for simpler workflows where
                  internal allocation is acceptable
                  
            \item All input arrays should be in \textbf{Fortran order} (column-major)
                  for optimal performance
                  
            \item Output arrays are \textbf{read-only} to maintain data integrity
          \end{itemize}
          
    \item \textbf{Complete Pipeline Functions}
          
          The Python interface provides two complete pipeline functions for processing
          trajectories with integrated outlier detection workflows.
          
          \begin{enumerate}
            \item \textbf{Process Trajectories with Per-Timepoint Percentiles}
                  
                  Processes multiple trajectories with percentile thresholds calculated separately
                  for each timepoint. This provides timepoint-specific sensitivity for
                  spike contribution analysis.
                  
                  \textbf{Arguments:}
                  \begin{itemize}
                    \item \texttt{trajectories}: 3D numpy array of shape \texttt{(n\_factors,
                          n\_samples, n\_timepoints)} in Fortran order
                    
                    \item \texttt{factor\_mask}: 1D boolean array of shape \texttt{(n\_factors,)}
                          indicating which factors to process
                          
                    \item \texttt{dependent\_idx}: 1-based index of the dependent variable
                          (will be excluded from processing)
                          
                    \item \texttt{mode}: Processing mode (1 = Normal, 2 = RAP)
                          
                    \item \texttt{percentile}: Percentile value for threshold calculation
                          (0.0--100.0)
                  \end{itemize}
                  \textbf{Returns:}
                  \begin{itemize}
                    \item Dictionary containing:
                          \begin{itemize}
                            \item \texttt{integrated\_contribs}: Integrated contributions \texttt{[n\_samples,
                                  n\_processed\_factors]}
                            
                            \item \texttt{spike\_contribs}: Spike contributions \texttt{[n\_timepoints,
                                  n\_samples, n\_processed\_factors]}
                            
                            \item \texttt{thresholds\_integrated}: Thresholds for integrated
                                  contributions \texttt{[n\_processed\_factors]}
                                  
                            \item \texttt{thresholds\_spike}: Thresholds for spike contributions
                                  \texttt{[n\_timepoints, n\_processed\_factors]}
                                  
                            \item \texttt{outliers\_integrated}: Outliers for integrated contributions
                                  \texttt{[n\_samples, n\_processed\_factors]} (boolean)
                                  
                            \item \texttt{outliers\_spike}: Outliers for spike contributions
                                  \texttt{[n\_timepoints, n\_samples, n\_processed\_factors]} (boolean)
                          \end{itemize}
                  \end{itemize}
                  \begin{lstlisting}[language=Python, caption={Complete Trajectory Processing with Per-Timepoint Thresholds}]
results = tox_process_trajectories(
    trajectories, factor_mask, dependent_idx, mode, percentile
)

# Access results
integrated_contribs = results['integrated_contribs']        
# [n_samples, n_processed]

spike_contribs = results['spike_contribs']                  
# [n_timepoints, n_samples, n_processed]

thresholds_integrated = results['thresholds_integrated']    
# [n_processed]

thresholds_spike = results['thresholds_spike']              
# [n_timepoints, n_processed]

outliers_integrated = results['outliers_integrated']        
# [n_samples, n_processed] (bool)

outliers_spike = results['outliers_spike']                  
# [n_timepoints, n_samples, n_processed] (bool)
                  \end{lstlisting}
                  
            \item \textbf{Process Trajectories with Global Percentile}
                  
                  Processes trajectories using a single global percentile threshold for all
                  timepoints in spike contributions. This provides consistent
                  sensitivity across all timepoints.
                  
                  \textbf{Arguments:}
                  \begin{itemize}
                    \item \texttt{trajectories}: 3D numpy array of shape \texttt{(n\_factors,
                          n\_samples, n\_timepoints)} in Fortran order
                    
                    \item \texttt{factor\_mask}: 1D boolean array of shape \texttt{(n\_factors,)}
                          indicating which factors to process
                          
                    \item \texttt{dependent\_idx}: 1-based index of the dependent variable
                          (will be excluded from processing)
                          
                    \item \texttt{mode}: Processing mode (1 = Normal, 2 = RAP)
                          
                    \item \texttt{percentile}: Percentile value for threshold calculation
                          (0.0--100.0)
                  \end{itemize}
                  \textbf{Returns:}
                  \begin{itemize}
                    \item Dictionary containing:
                          \begin{itemize}
                            \item \texttt{integrated\_contribs}: Integrated contributions \texttt{[n\_samples,
                                  n\_processed\_factors]}
                            
                            \item \texttt{spike\_contribs}: Spike contributions \texttt{[n\_timepoints,
                                  n\_samples, n\_processed\_factors]}
                            
                            \item \texttt{thresholds\_integrated}: Thresholds for integrated
                                  contributions \texttt{[n\_processed\_factors]}
                                  
                            \item \texttt{thresholds\_spike}: Thresholds for spike contributions
                                  \texttt{[n\_processed\_factors]} (1D array)
                                  
                            \item \texttt{outliers\_integrated}: Outliers for integrated contributions
                                  \texttt{[n\_samples, n\_processed\_factors]} (boolean)
                                  
                            \item \texttt{outliers\_spike}: Outliers for spike contributions
                                  \texttt{[n\_timepoints, n\_samples, n\_processed\_factors]} (boolean)
                          \end{itemize}
                  \end{itemize}
                  \begin{lstlisting}[language=Python, caption={Complete Trajectory Processing with Global Threshold}]
results = tox_process_trajectories_flat(
    trajectories, factor_mask, dependent_idx, mode, percentile
)

# Access results - note thresholds_spike is 1D in flat version
integrated_contribs = results['integrated_contribs']        
# [n_samples, n_processed]

spike_contribs = results['spike_contribs']                  
# [n_timepoints, n_samples, n_processed]

thresholds_integrated = results['thresholds_integrated']    
# [n_processed]

thresholds_spike = results['thresholds_spike']              
# [n_processed] (1D)

outliers_integrated = results['outliers_integrated']        
# [n_samples, n_processed] (bool)

outliers_spike = results['outliers_spike']                  
# [n_timepoints, n_samples, n_processed] (bool)
                  \end{lstlisting}
          \end{enumerate}
          
          \textbf{Key Differences Between Pipeline Functions:}
          \begin{itemize}
            \item \texttt{tox\_process\_trajectories}: Uses \textbf{per-timepoint
                  percentiles} for spike thresholds (\texttt{thresholds\_spike} shape: \texttt{[n\_timepoints,
                  n\_processed]})
            
            \item \texttt{tox\_process\_trajectories\_flat}: Uses \textbf{global
                  percentile} for spike thresholds (\texttt{thresholds\_spike} shape: \texttt{[n\_processed]})
          \end{itemize}
          
          \textbf{Pipeline Features:}
          \begin{itemize}
            \item \textbf{Automatic Factor Filtering}: Processes only factors marked
                  in \texttt{factor\_mask}, excluding the dependent variable
                  
            \item \textbf{Memory Efficient}: All output arrays use Fortran memory layout
                  and are marked as read-only
                  
            \item \textbf{Error Handling}: Comprehensive error checking with automatic
                  exception conversion
                  
            \item \textbf{Boolean Output}: Outlier masks are automatically converted
                  from C integers to Python boolean arrays
                  
            \item \textbf{Dimension Validation}: Input dimensions and factor mask consistency
                  are validated
          \end{itemize}
          
          \textbf{Usage Example:} 
          \begin{lstlisting}[language=Python, caption={Typical Pipeline Usage}]
import numpy as np

# Example data: 10 factors, 100 samples, 50 timepoints
trajectories = np.random.rand(10, 100, 50).astype(np.float64)
factor_mask = np.array([True, True, False, True, True, False, 
    True, True, True, False])
dependent_idx = 5  # 1-based index
mode = 1  # Normal mode
percentile = 95.0

# Process with per-timepoint thresholds
results = tox_process_trajectories(
    trajectories, factor_mask, dependent_idx, mode, percentile
)

# Process with global threshold  
results_flat = tox_process_trajectories_flat(
    trajectories, factor_mask, dependent_idx, mode, percentile
)
    \end{lstlisting}
    
    \item \textbf{Python Core Contribution Analysis Functions}
          
          \textbf{The Python interface provides comprehensive wrappers for
            trajectory contribution analysis, including EDF computation and various contribution
          calculation modes.}
          
          \begin{enumerate}
            \item \textbf{Empirical Distribution Function (EDF) Computation}
                  
                  Computes the empirical cumulative distribution function for observed data
                  values.
                  
                  \textbf{Arguments:}
                  \begin{itemize}
                    \item \texttt{values}: Array of observed data values (list or numpy
                          array)
                  \end{itemize}
                  \textbf{Returns:}
                  \begin{itemize}
                    \item Dictionary containing:
                          \begin{itemize}
                            \item \texttt{unique\_values}: Sorted unique data values (read-only
                                  numpy array)
                                  
                            \item \texttt{cdf\_values}: Corresponding cumulative frequencies
                                  between 0 and 1 (read-only)
                                  
                            \item \texttt{n\_unique}: Number of unique values found (int)
                          \end{itemize}
                  \end{itemize}
                  \begin{lstlisting}[language=Python, caption={Empirical Distribution Function Computation}]
result = compute_edf(values)

# Access results
unique_vals = result['unique_values']  
# Sorted unique values

cdf_vals = result['cdf_values']        
# Cumulative distribution values  

n_unique = result['n_unique']          
# Number of unique values
                  \end{lstlisting}
                  
            \item \textbf{Empirical Distribution Function (Expert Version)}
                  
                  Computes EDF using pre-sorted permutation for optimal performance.
                  
                  \textbf{Arguments:}
                  \begin{itemize}
                    \item \texttt{values}: Array of observed data values (list or numpy
                          array)
                          
                    \item \texttt{perm}: Pre-sorted permutation indices (1-based Fortran
                          style, must be sorted by values[perm])
                  \end{itemize}
                  \textbf{Returns:}
                  \begin{itemize}
                    \item Dictionary containing:
                          \begin{itemize}
                            \item \texttt{unique\_values}: Sorted unique data values (read-only
                                  numpy array)
                                  
                            \item \texttt{cdf\_values}: Corresponding cumulative frequencies
                                  between 0 and 1 (read-only)
                                  
                            \item \texttt{n\_unique}: Number of unique values found (int)
                          \end{itemize}
                  \end{itemize}
                  \begin{lstlisting}[language=Python, caption={Expert EDF Computation with Pre-sorted Permutation}]
result = compute_edf_expert(values, perm)

# Access results (same structure as compute_edf)
unique_vals = result['unique_values']
cdf_vals = result['cdf_values']
n_unique = result['n_unique']
                  \end{lstlisting}
          \end{enumerate}
\end{enumerate}